%-----------------------------------%
%----- Misc. Document Settings -----%
%-----------------------------------%
\documentclass[12pt,fleqn]{article}

\usepackage{preamble}
\usepackage{variables}

\begin{document}



%----- Sample Design -----%
\clearpage
\begin{table}[!htbp] 
    \begin{threeparttable}[b]
        \captionsetup{labelfont = bf, singlelinecheck = off, justification = raggedright, labelsep = none}   
        \caption{\\Sample Design}

        \small
        \begin{flushleft}
           {\raggedright
           \begin{tabular}{ l r}\\[-1.8ex]\hline \hline \\[-4.2ex]\\Firm-years listed in the Main Market of LSE during 2013 - 2023: &12,940\\ \quad Less firm-years missing Compustat matches &(5,140)\\ \quad Less firm-years missing variables from Compustat &(1,114)\\ \quad Less firm-years missing variables from IBES &(2,730)\\ \quad Less firm-years missing registration numbers &(263)\\ \quad Less firm-years before 2013-09-30 &\underline{(148)}\\ Firm-years provided to Gemini &3,545\\ \quad Less firm-years missing Companies House filings &(239)\\ \quad Less firm-years that do not report NGMs &(146)\\ \quad Less firm-years that do not report earnings-based NGMs &\underline{(278)}\\ Number of firm-years in sample &2,882\\ Average number of non-GAAP measures per firm-year &2.88\\Size of the firm-year-measure sample &8,287\\\\[-3.0ex]\hline \hline \\[-1.8ex]\end{tabular}
           \par}
        \end{flushleft}
        \label{table:Sample}

        \vspace{-6mm}
        \begin{tablenotes}
            \scriptsize
            \item \textbf{Note:} This table describes the screening process used to produce our sample.
        \end{tablenotes}

    \end{threeparttable}
\end{table}

%----- Descriptive Stats -----%
% \clearpage
\begin{table}[!htbp] 
    \begin{threeparttable}[b]
        \captionsetup{labelfont = bf, singlelinecheck = off, justification = raggedright, labelsep = none}   
        \caption{\\Descriptive Statistics}
        \small

        \begin{flushleft}        
        Panel A: Full sample    
        \end{flushleft}
        \vspace{-3mm}
        
        \begin{flushleft}
           {\raggedright
           {
\def\sym#1{\ifmmode^{#1}\else\(^{#1}\)\fi}
\begin{tabular}{l*{1}{ccccc}}
\\[-1.8ex]\hline \hline \\[-1.8ex]
                    &        mean&          sd&         p25&         p50&         p75\\
\midrule
\GAAPEPS        &       0.449&       0.876&       0.055&       0.241&       0.659\\
\STREETADJ   &       0.189&       0.667&       0.000&       0.065&       0.232\\
\MGMTADJ      &       0.387&       0.792&       0.042&       0.169&       0.477\\
\EXCESSADJ  &       0.197&       0.539&       0.000&       0.039&       0.230\\
\LossToProfit      &       0.138&       0.345&       0.000&       0.000&       0.000\\
\IncomeDecrease                 &       0.133&       0.339&       0.000&       0.000&       0.000\\
\AUDITED             &       0.766&       0.423&       1.000&       1.000&       1.000\\
\ONFACE &       0.364&       0.481&       0.000&       0.000&       1.000\\
\INNOTES            &       0.695&       0.461&       0.000&       1.000&       1.000\\
\NONGAAPISKAM     &       0.163&       0.370&       0.000&       0.000&       0.000\\
\MOSTPROM           &       0.210&       0.407&       0.000&       0.000&       0.000\\
\NUMNGMS              &       3.409&       1.258&       3.000&       3.000&       4.000\\
\NUMKAMS             &       3.616&       1.514&       3.000&       3.000&       4.000\\
\LOGASSETS           &      21.429&       1.775&      20.210&      21.219&      22.477\\
\LOSS                &       0.181&       0.385&       0.000&       0.000&       0.000\\
\BTM                 &       0.603&       0.563&       0.244&       0.445&       0.793\\
\REVGROWTH         &       0.032&       0.149&      -0.018&       0.026&       0.093\\
\VOLUME                 &       0.037&       0.044&       0.009&       0.021&       0.045\\
\ProcSub&       0.123&       0.328&       0.000&       0.000&       0.000\\
\ProcClass&       0.146&       0.354&       0.000&       0.000&       0.000\\
\ProcPol&       0.115&       0.319&       0.000&       0.000&       0.000\\
\ProcStds&       0.065&       0.247&       0.000&       0.000&       0.000\\
\ProcPast&       0.084&       0.277&       0.000&       0.000&       0.000\\
\ProcDisc&       0.107&       0.309&       0.000&       0.000&       0.000\\
\midrule
Observations        &        8287&            &            &            &            \\
\\[-1.8ex]\hline \hline \\[-1.8ex]
\end{tabular}
}

           \label{table:Descriptives}
           \par}
        \end{flushleft}

        \vspace{-6mm}
        \begin{flushleft}        
        Panel B: Split by Audited    
        \end{flushleft}
        \vspace{-3mm}
        
        \begin{flushleft}
           {\raggedright
           {
\def\sym#1{\ifmmode^{#1}\else\(^{#1}\)\fi}
\begin{tabular}{l*{3}{cc}}
\\[-1.8ex]\hline \hline \\[-1.8ex]
                    &\multicolumn{1}{c}{Audited = 0}&\multicolumn{1}{c}{Audited = 1}&\multicolumn{1}{c}{Difference}\\
\midrule
\GAAPEPS        &       0.478&       0.440&       0.038         \\
\STREETADJ   &       0.185&       0.190&      -0.006         \\
\MGMTADJ      &       0.518&       0.346&       0.172\sym{***}\\
\EXCESSADJ  &       0.327&       0.157&       0.170\sym{***}\\
\midrule
Observations        &        1939&        6348&        8287         \\
\\[-1.8ex]\hline \hline \\[-1.8ex]
\end{tabular}
}

           \label{table:Descriptives_by_audit}
           \par}
        \end{flushleft}

        \vspace{-6mm}
        \begin{tablenotes}
            \scriptsize
            \item \textbf{Note:} Panel A reports descriptive statistics for the full sample. Panel B reports descriptive statistics for selected key variables in subsamples split based on whether the non-GAAP measure is audited. *, **, and *** indicate p-values < .10, < .05, and < .01, respectively. Variable definitions are provided in Appendix A.
        \end{tablenotes}
    
    \end{threeparttable}
\end{table}



%----- Primary Hyp Test -----%
\clearpage
\begin{landscape}
\begin{table}[!htbp] 
    \begin{threeparttable}[b]
        \captionsetup{labelfont = bf, singlelinecheck = off, justification = raggedright, labelsep = none}   
        \caption{\\Are Audited Non-GAAP Adjustments Less Aggressive?}
        \centering
        \footnotesize

        \begin{flushleft}
           {\raggedright
           {
\def\sym#1{\ifmmode^{#1}\else\(^{#1}\)\fi}
\begin{tabular}{l*{5}{D{.}{.}{-1}}}
\\[-1.8ex]\hline \hline \\[-1.8ex]
                    &\multicolumn{1}{c}{(1)}&\multicolumn{1}{c}{(2)}&\multicolumn{1}{c}{(3)}&\multicolumn{1}{c}{(4)}&\multicolumn{1}{c}{(5)}\\
                    &\multicolumn{1}{c}{\shortstack{DV: \textit{\EXCESSADJ}}}&\multicolumn{1}{c}{\shortstack{DV: \textit{\EXCESSADJ}}}&\multicolumn{1}{c}{\shortstack{DV: \textit{\EXCESSADJ}}}&\multicolumn{1}{c}{\shortstack{DV: \textit{\EXCESSADJ}}}&\multicolumn{1}{c}{\shortstack{DV: \textit{\EXCESSADJ}}}\\
\midrule
\AUDITED             &      -0.170\sym{***}&      -0.158\sym{***}&      -0.159\sym{***}&      -0.219\sym{***}&      -0.299\sym{***}\\
                    &     (0.037)         &     (0.035)         &     (0.035)         &     (0.032)         &     (0.037)         \\
\GAAPEPS        &                     &       0.166\sym{***}&       0.166\sym{***}&       0.059\sym{***}&                     \\
                    &                     &     (0.026)         &     (0.026)         &     (0.020)         &                     \\
\MOSTPROM           &                     &       0.117\sym{***}&       0.116\sym{***}&       0.135\sym{***}&       0.144\sym{***}\\
                    &                     &     (0.025)         &     (0.025)         &     (0.023)         &     (0.024)         \\
\NONGAAPISKAM     &                     &      -0.004         &       0.004         &       0.018         &                     \\
                    &                     &     (0.034)         &     (0.038)         &     (0.028)         &                     \\
\NUMNGMS              &                     &       0.018\sym{***}&       0.016\sym{**} &       0.007         &                     \\
                    &                     &     (0.007)         &     (0.007)         &     (0.007)         &                     \\
\NUMKAMS             &                     &      -0.012         &      -0.010         &       0.009         &                     \\
                    &                     &     (0.012)         &     (0.014)         &     (0.007)         &                     \\
\LOGASSETS           &                     &       0.038\sym{***}&       0.037\sym{***}&      -0.006         &                     \\
                    &                     &     (0.011)         &     (0.012)         &     (0.040)         &                     \\
\LOSS                &                     &       0.134\sym{***}&       0.123\sym{***}&       0.047\sym{*}  &                     \\
                    &                     &     (0.034)         &     (0.033)         &     (0.024)         &                     \\
\BTM                 &                     &      -0.001         &      -0.002         &       0.072         &                     \\
                    &                     &     (0.028)         &     (0.029)         &     (0.047)         &                     \\
\REVGROWTH         &                     &       0.099\sym{*}  &       0.142\sym{***}&       0.097\sym{**} &                     \\
                    &                     &     (0.055)         &     (0.055)         &     (0.042)         &                     \\
\VOLUME                 &                     &       0.197         &       0.194         &       0.234         &                     \\
                    &                     &     (0.268)         &     (0.270)         &     (0.205)         &                     \\
Constant            &       0.327\sym{***}&      -0.649\sym{***}&      -0.632\sym{**} &       0.317         &       0.395\sym{***}\\
                    &     (0.034)         &     (0.238)         &     (0.251)         &     (0.842)         &     (0.029)         \\
\midrule
Fixed-effects       &        None         &        None         &         A,Y         &       F,A,Y         &         F*Y         \\
Observations        &       8,287         &       8,287         &       8,287         &       8,269         &       7,902         \\
Adjusted R-squared  &       0.018         &       0.112         &       0.117         &       0.414         &       0.432         \\
\\[-1.8ex]\hline \hline \\[-1.8ex]
\end{tabular}
}

           \par}
        \end{flushleft}

        \vspace{-6 mm}        
        \begin{tablenotes}
            \scriptsize
            \item Note: This table tests whether \textit{Mgmt. Excess Adj.} are smaller when the non-GAAP earnings are audited. We use OLS to estimate all the models. Column (1) includes no control variables or fixed effects. Column (2) adds firm-year–level and firm-year–measure–level control variables. Column (3) includes audit firm and year fixed effects. Column (4) includes client firm, audit firm, and year fixed effects. Column (5) includes client firm-by-year fixed effects and excludes firm-year–level controls because they do not vary within a given firm-year. All five columns use all firm-year–measures in the sample. The sample size in Column (4) decreases by 18, because those firms have only one-year observation in the sample. The sample size in Column (5) decreases by 385, because those firm-years have only one non-GAAP earnings measure. Standard errors are clustered at the firm-level. *, **, and *** indicate p-values of \textless\, .10, \textless\, .05, and \textless\, .01 respectively. Variable definitions are provided in Appendix A.
        \end{tablenotes}
        \label{table:street}

    \end{threeparttable}
\end{table}
\end{landscape}


%----- Alternative DVs -----%
\clearpage
\begin{table}[!htbp] 
    \begin{threeparttable}[b]
        \captionsetup{labelfont = bf, singlelinecheck = off, justification = raggedright, labelsep = none}   
        \caption{\\Alternative Dependent Variables}
        \small

        \begin{flushleft}
           {\raggedright
           {
\def\sym#1{\ifmmode^{#1}\else\(^{#1}\)\fi}
\begin{tabular}{l*{2}{D{.}{.}{-1}}}
\\[-1.8ex]\hline \hline \\[-1.8ex]
                    &\multicolumn{1}{c}{(1)}&\multicolumn{1}{c}{(2)}\\
                    &\multicolumn{1}{c}{\shortstack{Full Sample \\ DV: \textit{Negative Adj}}}&\multicolumn{1}{c}{\shortstack{\textit{Loss} = 1 \\ DV: \textit{Loss to Profit}}}\\
\midrule
\AUDITED             &       0.064\sym{***}&      -0.085\sym{***}\\
                    &     (0.011)         &     (0.028)         \\
\MOSTPROM           &      -0.083\sym{***}&       0.115\sym{***}\\
                    &     (0.009)         &     (0.024)         \\
Constant            &       0.095\sym{***}&       0.768\sym{***}\\
                    &     (0.008)         &     (0.022)         \\
\midrule
Fixed-effects       &         F*Y         &         F*Y         \\
Observations        &       7,902         &       1,428         \\
Adjusted R-squared  &       0.564         &       0.693         \\
\\[-1.8ex]\hline \hline \\[-1.8ex]
\end{tabular}
}

           \par}
        \end{flushleft}

        \vspace{-6 mm}
        
        \begin{tablenotes}
            \scriptsize
            \item \textbf{Note:} Column (1) tests whether the non-GAAP adjustments are more likely to be negative when the non-GAAP measure is audited. Column (2) tests whether the audited non-GAAP adjustments are less likely to convert a GAAP loss to a profit than the unaudited ones. Column (1) uses all firm-year-measures in the sample, except for the firm-years with only one non-GAAP earnings measure. Column (2) restricts the sample to the firm-years that report negative GAAP earnings. Both columns include client firm-by-year fixed effects. Standard errors are clustered at the firm-level. *, **, and *** indicate p-values of \textless\, .10, \textless\, .05, and \textless\, .01 respectively. Variable definitions are provided in Appendix A.
        \end{tablenotes}
        \label{table:negadj}
    \end{threeparttable}
\end{table}

%----- Face versus Notes -----%
\clearpage
% \begin{landscape}
\begin{table}[!htbp] 
    \begin{threeparttable}[b]
        \captionsetup{labelfont = bf, singlelinecheck = off, justification = raggedright, labelsep = none}   
        \caption{\\Face versus Notes}
        \small

        \begin{flushleft}
           {\raggedright
           {
\def\sym#1{\ifmmode^{#1}\else\(^{#1}\)\fi}
\begin{tabular}{l*{2}{D{.}{.}{-1}}}
\\[-1.8ex]\hline \hline \\[-1.8ex]
                    &\multicolumn{1}{c}{(1)}&\multicolumn{1}{c}{(2)}\\
                    &\multicolumn{1}{c}{\shortstack{\textit{\ONFACE} = 0 \\ DV: \textit{\EXCESSADJ}}}&\multicolumn{1}{c}{\shortstack{\textit{Full Sample} \\ DV: \textit{\EXCESSADJ}}}\\
\midrule
\AUDITED             &      -0.242\sym{***}&                     \\
                    &     (0.034)         &                     \\
\INNOTES            &                     &      -0.201\sym{***}\\
                    &                     &     (0.027)         \\
\ONFACE &                     &      -0.167\sym{***}\\
                    &                     &     (0.029)         \\
\MOSTPROM           &       0.157\sym{***}&       0.150\sym{***}\\
                    &     (0.038)         &     (0.026)         \\
Constant            &       0.343\sym{***}&       0.366\sym{***}\\
                    &     (0.022)         &     (0.024)         \\
\midrule
Fixed-effects       &         F*Y         &         F*Y         \\
Observations        &       4,309         &       7,902         \\
Adjusted R-squared  &       0.342         &       0.436         \\
\\[-1.8ex]\hline \hline \\[-1.8ex]
\end{tabular}
}

           \par}
        \end{flushleft}

        \vspace{-6mm}
        \begin{tablenotes}
            \scriptsize
            \item \textbf{Note:}  This table tests whether \textit{Mgmt. Excess Adj.} are smaller when the non-GAAP earnings are audited. Column (1) restricts the sample to the firm-year-measures that only appear in the MD\&A section of the annual report, or appear in both the MD\&A and the financial statements footnotes. Column (2) uses all firm-year-measures in the sample, except for the firm-years with only one non-GAAP earnings measure. Both columns include client firm-by-year fixed effects. Standard errors are clustered at the firm-level. *, **, and *** indicate p-values of \textless\, .10, \textless\, .05, and \textless\, .01 respectively. Variable definitions are provided in Appendix A.
        \end{tablenotes}
        \label{table:facevnootes}
    \end{threeparttable}
\end{table}
% \end{landscape}

%----- Specific Procedures: Panel A (portrait) -----%
\clearpage
\begin{table}[!htbp]
    \begin{threeparttable}[b]
        \captionsetup{labelfont = bf, singlelinecheck = off, justification = raggedright, labelsep = none}
        \caption{\\Specific Audit Procedures}
        \label{table:procs}

        \footnotesize
        \begin{flushleft}
           {\raggedright
           \footnotesize
\textit{Panel A: Procedure Descriptions and Examples}
\vspace{4pt}

\renewcommand{\arraystretch}{1.3}
\begin{tabular}{p{4.2cm} c p{9.8cm}}
\\[-1.8ex]\hline \hline \\[-1.8ex]
Procedure & Frequency & Example \\
\midrule
Challenge Classification & 90\% &
\textit{Kier Group plc, 30 June 2020} --- ``We reviewed the definition and classification of adjusting items in the Group's Annual Report, including the sub-categorisation of these items. In particular, we challenged whether it was appropriate to present certain costs within the Regional Southern Build business as restructuring and related charges, on the basis that they related to contract and tender positions.'' \\[6pt]
Substantive Testing & 75\% &
\textit{PZ Cussons plc, 31 May 2021} --- ``Substantively tested a sample of the adjusting items recognized in the year to supporting documentation, such as invoices or sale agreements and worked with tax specialists to assess the impact of the UK tax rate change.'' \\[6pt]
Ensure Aligns w/ Policy & 69\% &
\textit{Synthomer plc, 31 December 2023} --- ``We assessed the income and expenses classified as Special Items against the Group's accounting policies.'' \\[6pt]
Ensure Adequate Disclosure & 64\% &
\textit{Spirent Communications plc, 31 December 2023} --- ``Assessed whether the disclosures within the Group financial statements provided sufficient detail for the reader to understand the nature of these items and how adjusted results reconcile to statutory results.'' \\[6pt]
Ensure Consistent w/ Hist. & 52\% &
\textit{J D Wetherspoon plc, 29 July 2018} --- ``Ensuring that management's classification of exceptional items is consistent with prior year.'' \\[6pt]
Ensure Aligns w/ Standards & 39\% &
\textit{Mitie Group plc, 31 March 2017} --- ``We have benchmarked against market practice, including, but not limited to: the guidance published by the Financial Reporting Council in their thematic review; and the guidance included in the `Guidelines on Alternative Performance Measures', issued by the European Securities and Markets Authority (ESMA).'' \\
\\[-1.8ex]\hline \hline \\[-1.8ex]
\end{tabular}
\renewcommand{\arraystretch}{1.0}

           \par}
        \end{flushleft}

    \end{threeparttable}
\end{table}

%----- Specific Procedures: Panels B+C (landscape) -----%
\clearpage
\begin{landscape}
\begin{table}[!htbp]
    \addtocounter{table}{-1}
    \begin{threeparttable}[b]
        \captionsetup{labelfont = bf, singlelinecheck = off, justification = raggedright, labelsep = none}
        \caption{\\Specific Audit Procedures (continued)}

        \footnotesize
        \begin{flushleft}
           {\raggedright
           \footnotesize
\textit{Panel B: Criteria-based vs Non-criteria-based Procedures}
\vspace{4pt}

{
\def\sym#1{\ifmmode^{#1}\else\(^{#1}\)\fi}
\begin{tabular}{l*{1}{D{.}{.}{-1}}}
\\[-1.8ex]\hline \hline \\[-1.8ex]
                    &\multicolumn{1}{c}{(1)}\\
                    &\multicolumn{1}{c}{DV: \textit{\EXCESSADJ}}\\
\midrule
\AuditNoBenchGroup&      -0.213\sym{***}\\
                    &     (0.032)         \\
\AuditAndBenchGroup&      -0.293\sym{***}\\
                    &     (0.046)         \\
\midrule
Fixed-effects       &       F,A,Y         \\
Controls            &         Yes         \\
Observations        &       8,269         \\
Adjusted R-squared  &       0.414         \\
p-value for Difference&       0.033         \\
\\[-1.8ex]\hline \hline \\[-1.8ex]
\end{tabular}
}

\vspace{6pt}

\textit{Panel C: Individual Procedures}
\vspace{4pt}

{
\def\sym#1{\ifmmode^{#1}\else\(^{#1}\)\fi}
\begin{tabular}{l*{6}{D{.}{.}{-1}}}
\\[-1.8ex]\hline \hline \\[-1.8ex]
                    &\multicolumn{3}{c}{Non-criteria-based}                           &\multicolumn{3}{c}{Criteria-based}                               \\\cmidrule(lr){2-4}\cmidrule(lr){5-7}
                    &\multicolumn{1}{c}{(1)}&\multicolumn{1}{c}{(2)}&\multicolumn{1}{c}{(3)}&\multicolumn{1}{c}{(4)}&\multicolumn{1}{c}{(5)}&\multicolumn{1}{c}{(6)}\\
                    &\multicolumn{1}{c}{\shortstack{Challenge\\Classif.}}&\multicolumn{1}{c}{\shortstack{Substantive\\Testing}}&\multicolumn{1}{c}{\shortstack{Adequate\\Disclosure}}&\multicolumn{1}{c}{\shortstack{Aligns w/\\Policy}}&\multicolumn{1}{c}{\shortstack{Consistent\\w/ Hist.}}&\multicolumn{1}{c}{\shortstack{Aligns w/\\Standards}}\\
\midrule
\AuditNoProc       &      -0.212\sym{***}&      -0.219\sym{***}&      -0.218\sym{***}&      -0.216\sym{***}&      -0.210\sym{***}&      -0.217\sym{***}\\
                    &     (0.031)         &     (0.031)         &     (0.031)         &     (0.033)         &     (0.032)         &     (0.031)         \\
\AuditAndProc      &      -0.262\sym{***}&      -0.224\sym{***}&      -0.227\sym{***}&      -0.240\sym{***}&      -0.309\sym{***}&      -0.258\sym{***}\\
                    &     (0.058)         &     (0.059)         &     (0.063)         &     (0.045)         &     (0.052)         &     (0.073)         \\
\midrule
Fixed-effects       &       F,A,Y         &       F,A,Y         &       F,A,Y         &       F,A,Y         &       F,A,Y         &       F,A,Y         \\
Controls            &         Yes         &         Yes         &         Yes         &         Yes         &         Yes         &         Yes         \\
Observations        &       8,269         &       8,269         &       8,269         &       8,269         &       8,269         &       8,269         \\
Adjusted R-squared  &       0.414         &       0.414         &       0.414         &       0.414         &       0.415         &       0.414         \\
p-value for Difference&       0.285         &       0.909         &       0.869         &       0.538         &       0.018         &       0.503         \\
\\[-1.8ex]\hline \hline \\[-1.8ex]
\end{tabular}
}

           \par}
        \end{flushleft}
        \vspace{-6mm}

        \begin{tablenotes}
            \scriptsize
            \vspace{3mm}
            \item \textbf{Note:} This table reports results from OLS regressions of \textit{Mgmt. Excess Adj.} on audit procedure indicators. Panel A describes each procedure and its frequency. Panel B partitions audited observations into criteria-based (\textit{Ensure Aligns w/ Policy}, \textit{Ensure Consistent w/ Hist.}, and \textit{Ensure Aligns w/ Standards}) versus non-criteria-based (\textit{Challenge Classification}, \textit{Substantive Testing}, and \textit{Ensure Adequate Disclosure}) procedure groups using indicators based on procedure counts, with ties activating both. Panel C reports separate regressions for each individual procedure. All regressions include firm, auditor, and year fixed effects and the full set of control variables. Standard errors are clustered by firm and reported in parentheses. *, **, and *** denote significance at the 10\%, 5\%, and 1\% levels, respectively.
        \end{tablenotes}

    \end{threeparttable}
\end{table}
\end{landscape}

%----- Value Relevance -----%
\clearpage
\begin{table}[!htbp] 
    \begin{threeparttable}[b]
        \captionsetup{labelfont = bf, singlelinecheck = off, justification = raggedright, labelsep = none}   
        \caption{\\Are Audited Non-GAAP Earnings More Value Relevant?}
        \small

        \begin{flushleft}
           {\raggedright
           {
\def\sym#1{\ifmmode^{#1}\else\(^{#1}\)\fi}
\begin{tabular}{l*{2}{D{.}{.}{-1}}}
\\[-1.8ex]\hline \hline \\[-1.8ex]
                    &\multicolumn{1}{c}{(1)}&\multicolumn{1}{c}{(2)}\\
                    &\multicolumn{1}{c}{\shortstack{DV: \textit{\PRICETHIRTY}}}&\multicolumn{1}{c}{\shortstack{DV: \textit{\PRICETHIRTY}}}\\
\midrule
\EXCESSADJ  &       1.003\sym{***}&       0.021         \\
                    &     (0.316)         &     (0.384)         \\
\AUDITED             &                     &       0.058         \\
                    &                     &     (0.189)         \\
\ExcessXAudited&                     &       1.629\sym{***}\\
                    &                     &     (0.541)         \\
\STREETADJ   &       2.506\sym{***}&       2.567\sym{***}\\
                    &     (0.471)         &     (0.458)         \\
\GAAPEPS        &       4.064\sym{***}&       4.105\sym{***}\\
                    &     (0.487)         &     (0.474)         \\
\BV           &       0.393         &       0.389         \\
                    &     (0.270)         &     (0.265)         \\
Constant            &       6.534\sym{***}&       6.474\sym{***}\\
                    &     (1.161)         &     (1.138)         \\
\midrule
Fixed-effects       &       F,A,Y         &       F,A,Y         \\
Observations        &       8,269         &       8,269         \\
Adjusted R-squared  &       0.901         &       0.901         \\
\\[-1.8ex]\hline \hline \\[-1.8ex]
\end{tabular}
}

           \par}
        \end{flushleft}
        \vspace{-6mm}
        
        \begin{tablenotes}
            \scriptsize
            \item \textbf{Note:} This table examines the value relevance of the non-GAAP adjustments. Column (1) does not include any variables capturing the audited status of the non-GAAP earnings measure. Column (2) includes both the \textit{Audited} variable and its interaction with \textit{Mgmt. Excess Adj.} Both columns use all firm-year-measures, except for the firms that have only one-year observation in the sample. Both columns include client firm, audit firm, and year fixed effects. Robust standard errors are clustered at the firm-level. *, **, and *** indicate p-values of \textless\, .10, \textless\, .05, and \textless\, .01 respectively. Variable definitions are provided in Appendix A.
        \end{tablenotes}
        \label{table:valrel}
    \end{threeparttable}
\end{table}

\end{document}
